\section{Riemann surfaces -- the symplectic case}\label{section3}
The conformal invariance of the Dirac equation in two dimensions has a more important consequence since it can be defined in terms of the holomorphic structure and so becomes part of the algebraic geometry of a compact projective curve $C$ of genus $g$. A spin structure is defined by the holomorphic structure on a square root $K^{1/2}$ of the canonical bundle $K$ and the Dirac operator is just $\slashed D=\bar\partial\colon C^{\infty}\big(K^{1/2}\big)\rightarrow C^{\infty}\big(K^{1/2}\bar K\big)$.
In this case, the index theorem, which is the Riemann--Roch theorem, gives $\dim H^0\big(C,K^{1/2}\big)-\dim H^1\big(C, K^{1/2}\big)=0$.
When coupled to a vector bundle $E$ with connection, the $(0,1)$ component of the connection defines a $\bar\partial$-operator on $E$ and the null space of the coupled Dirac operator is $H^0\big(C,E\otimes K^{1/2}\big)$ and if $c_1(E)=0$ the index is again zero.

Suppose we are given a holomorphic principal $G$-bundle over $C$ and we take a symplectic representation to yield a vector bundle $E$. The moment map for a symplectic action of $G$ on a vector space $V$ is a quadratic map from $V$ to $\lie{g}^*$, the dual of the Lie algebra. Then, using a~bi-invariant metric, a holomorphic section of $E\otimes K^{1/2}$ gives a section of $\lie{g}\otimes K$. These are called Higgs fields and the moduli space of bundles together with Higgs fields has gained importance in various fields. The construction above of special Higgs fields was suggested by the physicist D.~Gaiotto \cite{Gaiotto} and described in~\cite{Hitchin2}.

We shall consider the irreducible symplectic representations of ${\rm SL}(2,\C)$. These are the symmetric powers $S^m\C^2$ for odd $m$ of the basic representation on $\C^2$ or equivalently the polynomials in one variable of odd degree $m$ under the action
$p(z)\mapsto (cz+d)^mp((az+b)/(cz+d)).$
 Set $m=2k-1$ and write $p(z)=a_0z^m+a_1z^{m-1}+\dots+a_m$, then the invariant symplectic form on the space $S^m\C^2$ is up to a constant
\begin{equation}
\omega=\sum_{\ell=0}^{k-1}(-1)^{\ell}\ell!(m-\ell)!{\rm d}a_i\wedge {\rm d}a_{m-i}.\label{symp}
\end{equation}
The even symmetric powers $S^{m}\C^2$ are the orthogonal irreducible representations and $m=2$ is the adjoint representation. The moment map can be viewed as $\mu\colon S^m\C^2\otimes S^m\C^2\rightarrow S^2\C^2$ given by contraction using the symplectic form.

 First consider $m=1$.
We start with a rank 2 holomorphic vector bundle $V$ with $\Lambda^2V$ trivial. Multiples of a holomorphic section $\psi$ of $V\otimes K^{1/2}$ generate a subbundle $L^*\subset V$ with quotient~$L$. Then $\psi$ is a section of the line bundle $L^*K^{1/2}$. This has degree $g-1-\deg L$. So to construct such a~$\psi$ we take an effective divisor $D$ of degree $< g-1$ and define $L$ to be $K^{1/2}(-D)$ and the vector bundle to be given by an extension class in $H^1\big(C,L^{-2}\big)=H^1(C,K(-2D))$.

The moment map in this case gives an everywhere nilpotent Higgs field $\Phi=\psi\otimes \psi$, but for $m>1$ this is not necessarily so and we can make use of the relation to Higgs bundles and adopt the spectral curve approach.
We start with a holomorphic section $q$ of $K^2$ with simple zeros and define the curve $S$ in the total space $\vert K\vert$ of $K$ (i.e., the cotangent bundle of $C$) by the equation $x^2=\pi^*q$. Here $\pi\colon \vert K\vert \rightarrow C$ is the projection and $x$ is the tautological section of $\pi^*K$ on the total space $\vert K\vert $. Then $S$ is a smooth curve of genus $4g-3$ and has a covering involution $\sigma$ defined by $x\mapsto-x$. The direct image of a line bundle $L$ on $S$ is a rank 2 bundle $V$ on $C$ and if $L\cong U\pi^*K^{1/2}$ and $\sigma^*U\cong U^*$ then $\Lambda^2V$ is trivial (see for example~\cite{BeauvilleA/NarasimhanM/RamananS}). The condition on $U$ can be interpreted as saying it lies in the Prym variety of the map $\pi\colon S\rightarrow C$, an abelian variety of dimension $4g-3-g=3g-3$.

Proposition~3 of \cite{Hitchin2} shows how to find a solution to the Dirac operator coupled to $S^mV$ from a line bundle of order $m$ on $S$ -- a point of finite order in the Prym variety. For simplicity we consider the value $m=3$. For our purposes we are not interested in the Higgs field, only its spectral curve.

\begin{proposition} Let $U$ be a non-trivial line bundle on $S$ such that $U^3$ is trivial. Then the direct image $V$ of $U\pi^*K^{1/2}$ admits a holomorphic section of $S^3V\otimes K^{1/2}$.
\end{proposition}

\begin{proof} The definition of direct image $V=\pi_*L$ is that for each open set $W\subset C$, $H^0(W, \pi_*L)=H^0\big(\pi^{-1}(W), L\big)$ so that evaluation of a local section of $L$ expresses the pull-back $\pi^*V$ as an extension{\samepage
\begin{equation}
0\rightarrow L^*\rightarrow \pi^*V\rightarrow L\rightarrow 0.
\label{V1}
\end{equation}
and this induces a filtration of $\pi^*S^3V$ into subbundles $V_0\subset V_1\subset V_2\subset V_3$ with $V_0\cong L^{-3}$.}

Consider the subbundle $V_1$. Comparing with the filtration by degree of polynomials $a_0z^3+a_1z^{2}+a_2z+a_3$ it is the subspace $a_0=a_{1}=0$, the polynomials which vanish to order $2$ at $\infty$. It is maximally isotropic with respect to the symplectic form \eqref{symp} with $m=3$.
It follows that $\pi^*S^3V$ is an extension
\begin{equation}
0\rightarrow V_1\rightarrow \pi^*S^3V\rightarrow V_1^*\rightarrow 0
\label{ext}
\end{equation}

Now $\pi^*V$ is invariant by $\sigma$ so from \eqref{V1} we have another extension
\begin{displaymath}0\rightarrow \sigma^*L^*\rightarrow \pi^*V\rightarrow \sigma^*L\rightarrow 0\end{displaymath}
and this splits \eqref{V1} outside the divisor given by $x=0$. The extension class is then supported on $x=0$, the fixed point set of $\sigma$, where $L\cong \sigma^*L$. In Dolbeault terms
we have a section of $L\sigma^*L^*$ on $x=0$, extend it holomorphically to a neighbourhood and then everywhere as a $C^{\infty}$ section $s$ supported in a neighbourhood of the fixed point set. The extension class in $H^1\big(S,L^{-2}\big)$ is then represented by $\bar\partial s/x\in \Omega^{0,1}\big(S, L^{-2}\big)$.

The Dolbeault representative for the holomorphic structure of $V$ as an extension gives a $C^{\infty}$-splitting $\pi^*V=L\oplus L^*$ and a $\bar\partial $-operator $\bar\partial +\alpha$, where $\alpha\in \Omega^{0,1}\big(S, L^{-2}\big)$. Then $\pi^*S^3V=L^{3}\oplus L\oplus L^{-1}\oplus L^{-3}$ and a holomorphic section of $\pi^*\big(S^3V\otimes K^{1/2}\big)$ can be written as $a_0z^3+a_1z^{2}+a_1z+a_0$, where
\begin{displaymath}\bar\partial a_0=0,\qquad \bar\partial a_1+3\alpha a_0=0,\qquad \bar\partial a_2+2\alpha a_1=0, \qquad \bar\partial a_3+\alpha a_2=0\end{displaymath}
and $a_0\in \Omega^0\big(S, L^3\pi^*K^{1/2}\big)$, $a_1\in \Omega^0\big(S, L\pi^*K^{1/2}\big)$, etc.

The line bundle $U^3$ is trivial. Let $u$ be a trivialization such that $\sigma^*u=u^{-1}$ and set $a_0=ux^2$. Then
$3\alpha a_0=3u(\bar\partial s) x=\bar\partial (3usx)$
so we can take $a_1=-3usx$ and extend $a_0$ to get a holomorphic section $v$ of $V_1^*\otimes \pi^*K^{1/2}$. This vanishes on $x=0$ because $a_0$ and $a_1$ do. Any two choices of $a_1$ differ by an element of $H^0\big(S, L\pi^*K^{1/2}\big)$ and if they are divisible by $x$ then this is equivalent to sections of $L\pi^*K^{-1/2}=U$. Thus if $U$ is non-trivial there is a unique extension divisible by $x$.

To extend $v$ to a section of $\pi^*\big(S^3V\otimes K^{1/2}\big)$ we observe that the vector bundle extension~\eqref{ext} is defined by a class in $H^1\big(S, \Hom(V_1^*,V_1)\big)$ also of the form $\bar\partial a/x$ and since $v$ as constructed is divisible by~$x$,~$v$ extends. In this case any two extensions differ by $v'\in H^0\big(S,V_1\otimes \pi^*K^{1/2}\big)$.

But we have
 \begin{displaymath}0\rightarrow L^{-3}=V_{0}\rightarrow V_{1}\rightarrow L^{-1}\rightarrow 0\end{displaymath}
and $L^{-3}\cong \pi^*K^{-3/2}$, so $V_0\otimes \pi^*K^{1/2}$ is a line bundle of negative degree and so has no sections. Moreover $L^{-1}\pi^*K^{1/2} \cong U^{-1}$ has degree zero and is non-trivial so this has no holomorphic sections either. It follows from the exact cohomology sequence that $H^0\big(S, V_1\otimes \pi^*K^{1/2}\big)=0$ and there is a unique extension.

 This gives us a section $v''$ of $\pi^*\big(S^3V\otimes K^{1/2}\big)$ on $S$. But our choice of trivialization $u$ implies that $\sigma$ takes $ux^2$ to $u^{-1}x^2$ and by uniqueness $\sigma^*v''=v''$. Since $\pi^*V$ is pulled back from~$C$, $v''$~descends to a section $\psi$ of $S^3V\otimes K^{1/2}$.
\end{proof}
